% Extracción quirúrgica del propietario V2 05d: líneas 286--512.
% El resto permanece en este archivo como procedencia, pero no se publica.
\iffalse
% Macroestructura de coherencia estrictamente interna del continuo discreto.
% No usa geometrías, clases ni conclusiones clásicas.

\chapter{Macroestructura interna de celdas, coherencia y completitud}
\label{chap:pdfv2-coh-interna}

\begin{statusbox}
La única entrada de este capítulo es la estructura discreta del continuo. La
rama \(\mathrm{coh}\) se construye sobre los anillos
\(A_N=\mathbb Z/3^N\mathbb Z\), las celdas APP, sus historias y sus
terminales. Su cadena no ablativa es
\[
 \mathcal C_{\rm cont}^{\rm disc}
 \supset\mathcal C_{\rm cont}^{\rm coh}
 \xrightarrow{\operatorname{Res}_{\rm coh}}
 \mathcal G_{\rm coh}
 \xrightarrow{\operatorname{pr}_{X,{\rm coh}}}
 X_{\chi,{\rm coh}}
 \xrightarrow{\mathcal A_{\rm coh}}
 \mathfrak M_{\rm coh}
 \xrightarrow{\mathcal P_{\rm coh}}
 \mathcal C_{\rm coh}^{\rm HMT}.
\]
No aparece como entrada ni como salida ningún objeto clásico. Cada imagen
conserva literalmente su historia generadora.
\end{statusbox}

\section{Dominio tipado y separación de índices}
\label{sec:pdfv2-coh-domain}

Para no identificar profundidad con precisión, se usa el conjunto dirigido
\begin{equation}\label{eq:pdfv2-coh-lambda}
 \Lambda:=\mathbb N_{\rm prof}\times\mathbb N_{\rm prec},
 \qquad
 \lambda=(d,N),
 \qquad
 A_N:=\mathbb Z/3^N\mathbb Z.
\end{equation}
Si \(\lambda\leq\lambda'=(d',N')\), la transición
\begin{equation}\label{eq:pdfv2-coh-rho-lambda}
 \rho_{\lambda',\lambda}:
 \mathcal C_{{\rm cont},\lambda'}^{\rm disc}
 \longrightarrow
 \mathcal C_{{\rm cont},\lambda}^{\rm disc}
\end{equation}
trunca la historia después de la profundidad \(d\) y reduce coeficientes por
\(A_{N'}\twoheadrightarrow A_N\), conservando hoja, orientación, residuo,
cociente, carry, memoria, frontera, ruta, supervivencia y terminal.

La subestructura \(\mathcal C_{\rm cont}^{\rm coh}\) consta de las historias
\(x\in\mathcal C_{\rm cont}^{\rm disc}\) que cumplen, en todo nivel:
\begin{enumerate}
 \item el prefijo de \(x\) porta las celdas orientadas y basadas definidas en
       la sección siguiente;
 \item cada refinamiento induce un mapa del complejo de la celda;
 \item truncación, reducción de coeficientes, orientación y extensión
       nonádica conmutan;
 \item las multisecciones de extensión superviviente son no vacías a toda
       profundidad;
 \item el retorno nonádico conserva la diferencia terminal entre la historia
       y su continuación.
\end{enumerate}
Estas condiciones definen el dominio de la rama. No se atribuyen a una
historia arbitraria sin verificarlas.

\section{Celda APP sobre \texorpdfstring{\(A_N\)}{AN}}
\label{sec:pdfv2-coh-cell}

Una celda orientada es
\[
 \nu=(h_\nu,I_\nu,J_\nu,\epsilon_\nu),
\]
donde \(h_\nu\) es una historia finita, \(I_\nu,J_\nu\) son conjuntos
ordenados no vacíos de puertos,
\[
 a_\nu:=|I_\nu|\geq1,
 \qquad
 b_\nu:=|J_\nu|\geq1,
\]
y \(\epsilon_\nu\in\{+1,-1\}\) es su orientación.

Se define
\begin{equation}\label{eq:pdfv2-coh-kappa}
 \begin{aligned}
 \kappa_{\nu,N}:
 A_N^{I_\nu}\oplus A_N^{J_\nu}
 &\longrightarrow A_N^{I_\nu\times J_\nu},\\
 \kappa_{\nu,N}(u,v)_{ij}&:=u_i-v_j.
 \end{aligned}
\end{equation}
Elegidos los puertos base \(i_0\in I_\nu\), \(j_0\in J_\nu\), sea
\begin{equation}\label{eq:pdfv2-coh-delta}
 \begin{aligned}
 \delta_{\nu,N}:A_N^{I_\nu\times J_\nu}
 &\longrightarrow
 A_N^{(I_\nu\setminus i_0)\times(J_\nu\setminus j_0)},\\
 \delta_{\nu,N}(w)_{ij}
 &:={}
 w_{ij}-w_{i j_0}-w_{i_0j}+w_{i_0j_0}.
 \end{aligned}
\end{equation}

\begin{proposition}[Sucesión exacta de la celda]
\label{prop:pdfv2-coh-cell-exact}
Para toda celda \(\nu\) y toda precisión \(N\), la sucesión
\begin{equation}\label{eq:pdfv2-coh-exact-sequence}
 0\longrightarrow A_N
 \xrightarrow{\ \Delta\ }
 A_N^{I_\nu}\oplus A_N^{J_\nu}
 \xrightarrow{\ \kappa_{\nu,N}\ }
 A_N^{I_\nu\times J_\nu}
 \xrightarrow{\ \delta_{\nu,N}\ }
 A_N^{(a_\nu-1)(b_\nu-1)}
 \longrightarrow0,
\end{equation}
donde
\[
 \Delta(t)=\bigl((t)_{i\in I_\nu},(t)_{j\in J_\nu}\bigr),
\]
es exacta y escindida sobre \(A_N\).
\end{proposition}

\begin{proof}
Si \(\kappa_{\nu,N}(u,v)=0\), entonces \(u_i=v_j\) para todo \(i,j\),
por lo que \((u,v)\) es diagonal. Así,
\(\ker\kappa_{\nu,N}=\operatorname{im}\Delta\). La sustitución directa da
\(\delta_{\nu,N}\kappa_{\nu,N}=0\).

Si \(\delta_{\nu,N}(w)=0\), póngase
\[
 u_i=w_{i j_0},
 \qquad
 v_j=w_{i_0j_0}-w_{i_0j}.
\]
Entonces
\[
 u_i-v_j=w_{i j_0}+w_{i_0j}-w_{i_0j_0}=w_{ij},
\]
de modo que \(\ker\delta_{\nu,N}=\operatorname{im}\kappa_{\nu,N}\). Para
probar la sobreyectividad de \(\delta_{\nu,N}\), se anulan la fila y la
columna base y se colocan arbitrariamente las coordenadas deseadas en las
posiciones restantes. Esas mismas fórmulas proporcionan escisiones.
\end{proof}

La celda compleja, en grados \(1\to0\), es el cono homológico de
\(\kappa_{\nu,N}\) con la convención que coloca su fuente en grado \(1\):
\begin{equation}\label{eq:pdfv2-coh-cell-complex}
 C_{\nu,N}:=\operatorname{Cone}(\kappa_{\nu,N})
 :=\left[
 A_N^{I_\nu}\oplus A_N^{J_\nu}
 \xrightarrow{\ \kappa_{\nu,N}\ }
 A_N^{I_\nu\times J_\nu}
 \right].
\end{equation}
La proposición anterior da
\begin{equation}\label{eq:pdfv2-coh-cell-homology}
 H_1(C_{\nu,N})\simeq A_N,
 \qquad
 H_0(C_{\nu,N})\simeq A_N^{(a_\nu-1)(b_\nu-1)}.
\end{equation}
El defecto de celda queda así construido, no añadido como un cardinal sin
mapa.

\section{Orientación y carry de la celda}
\label{sec:pdfv2-coh-orientation-carry}

Sean
\[
 \tau_1(u,v)=(v,u),
 \qquad
 P(w)_{ji}=w_{ij}.
\]
El intercambio orientado es el mapa de complejos
\begin{equation}\label{eq:pdfv2-coh-theta}
 \Theta_{\nu,N}:=(\tau_1,-P):
 C_{a_\nu,b_\nu,N}\longrightarrow C_{b_\nu,a_\nu,N},
\end{equation}
porque
\begin{equation}\label{eq:pdfv2-coh-orientation}
 \kappa_{b_\nu,a_\nu,N}\tau_1
 =-P\kappa_{a_\nu,b_\nu,N}.
\end{equation}
Además,
\(\Theta_{\nu^{\rm op},N}\Theta_{\nu,N}=I\). La ruta opuesta no se
identifica con la original: cambia el signo de grado cero y conserva el acto
de inversión.

Para \(r\in A_N\), sea
\(\langle r\rangle_N\in\{0,\ldots,3^N-1\}\) su representante canónico. El
borrow/carry local es
\begin{equation}\label{eq:pdfv2-coh-beta}
 \beta_N(r,s):=
 \frac{\langle r-s\rangle_N-\langle r\rangle_N+\langle s\rangle_N}
 {3^N}\in\{0,1\}.
\end{equation}
Así,
\begin{equation}\label{eq:pdfv2-coh-carry-cell}
 \langle u_i-v_j\rangle_N
 =\langle u_i\rangle_N-\langle v_j\rangle_N
  +3^N\beta_N(u_i,v_j).
\end{equation}
La matriz
\[
 \operatorname{Carr}_{\nu,N}(u,v)
 :=\bigl(\beta_N(u_i,v_j)\bigr)_{ij}
\]
se conserva junto a \(\kappa_{\nu,N}(u,v)\): residuo y carry son dos
coordenadas distintas.

Cuando los tamaños de puertos proceden de las dos hojas APP, la admisibilidad
de la celda incluye la identidad entera
\begin{equation}\label{eq:pdfv2-coh-app-defect}
 (a_\nu-1)(b_\nu-1)
 =\rho_{\Pi,\nu}-\rho_{\Sigma,\nu}+1
  +9(c_{\Pi,\nu}-c_{\Sigma,\nu}).
\end{equation}
La igualdad se exige únicamente a una celda con esa genealogía; no se afirma
para enteros desprovistos de hojas APP.

Si \(L_N(u)\) es la elevación entera canónica de la matriz de una flecha
\(u\), se define el carry de composición por
\begin{equation}\label{eq:pdfv2-coh-carry-composition}
 L_N(v)L_N(u)=L_N(vu)+3^N c_N(v,u).
\end{equation}
La asociatividad de tres flechas da el cociclo exacto
\begin{equation}\label{eq:pdfv2-coh-carry-cocycle}
 c_N(w,vu)+L_N(w)c_N(v,u)
 =c_N(wv,u)+c_N(w,v)L_N(u).
\end{equation}

\section{Categoría de historias y functor celular}
\label{sec:pdfv2-coh-category}

Para un prefijo \(x_{\leq d}\), sea \(J_d(x)\) la categoría finita cuyos
objetos son todas sus celdas \(\nu\). Sus flechas están generadas por:
\begin{enumerate}
 \item refinamientos \(u:\nu\to\nu'\) con aplicaciones sobreyectivas de
       puertos
       \[
        \alpha_u:I_{\nu'}\twoheadrightarrow I_\nu,
        \qquad
        \beta_u:J_{\nu'}\twoheadrightarrow J_\nu;
       \]
 \item los intercambios de orientación \(\Theta_\nu\);
 \item las extensiones nonádicas \(\Gamma_9\nu\);
 \item identidad y composición, sometidas a las relaciones de frontera,
       ruta, orientación y carry de la historia.
\end{enumerate}
La finitud se refiere al prefijo; el sistema de todos los prefijos no se
declara finito.

El functor celular es
\begin{equation}\label{eq:pdfv2-coh-gamma-functor}
 \Gamma_{d,N}:J_d(x)\longrightarrow
 \operatorname{Ch}^{[0,1]}(A_N),
 \qquad
 \Gamma_{d,N}(\nu)=C_{\nu,N}.
\end{equation}
Para una flecha de refinamiento \(u:\nu\to\nu'\), actúa mediante
\begin{align}
 \Gamma_{d,N}(u)_1(p,q)
 &:=(p\circ\alpha_u,q\circ\beta_u),
 \label{eq:pdfv2-coh-gamma-one}\\
 \Gamma_{d,N}(u)_0(w)
 &:=w\circ(\alpha_u\times\beta_u).
 \label{eq:pdfv2-coh-gamma-zero}
\end{align}
Se verifica coordenada a coordenada
\begin{equation}\label{eq:pdfv2-coh-gamma-chain-map}
 \kappa_{\nu',N}\Gamma_{d,N}(u)_1
 =\Gamma_{d,N}(u)_0\kappa_{\nu,N}.
\end{equation}
Sobre una inversión de orientación actúa por
\(\Theta_{\nu,N}\), y sobre composiciones actúa por composición, guardando
separadamente \(c_N(v,u)\).

La reducción \(r_{N+1,N}:A_{N+1}\twoheadrightarrow A_N\) satisface
\begin{equation}\label{eq:pdfv2-coh-kappa-reduction}
 r_{N+1,N}^{I\times J}\kappa_{\nu,N+1}
 =\kappa_{\nu,N}
  \bigl(r_{N+1,N}^{I}\oplus r_{N+1,N}^{J}\bigr).
\end{equation}
Por tanto, el diagrama de celdas es natural simultáneamente en historia y
precisión.

\fi
\section{Módulo de rutas, totalización celular y terminal nonádico}
\label{sec:pdfv2-coh-route-module}

Sea \(\operatorname{Term}_d(x)\) el conjunto de terminales de memoria del
prefijo. Para \(\nu\in J_d(x)\), defínase
\begin{equation}\label{eq:pdfv2-coh-w-module}
 W_{d,N}(\nu):=
 A_N\bigl[
 \operatorname{Path}_{J_d(x)}(\nu,\operatorname{Term}_d(x))
 \bigr].
\end{equation}
Es el módulo libre sobre todas las rutas terminales; no selecciona una. Si
\(u:\nu\to\nu'\), la acción contravariante es
\begin{equation}\label{eq:pdfv2-coh-w-action}
 W_{d,N}(u):W_{d,N}(\nu')\longrightarrow W_{d,N}(\nu),
 \qquad
 [\lambda]\longmapsto[\lambda\circ u].
\end{equation}
Cada generador conserva orientación, carry, frontera, ruta y memoria.

\section{Complejo bar bilateral y totalización de la bicadena}
\label{sec:pdfv2-coh-bar}

Para \(q\geq0\) y \(r\in\{0,1\}\), sea
\begin{equation}\label{eq:pdfv2-coh-bar-bigraded}
 B_{q,r}^{d,N}(x):=
 \bigoplus_{\nu_0\xrightarrow{u_1}\cdots\xrightarrow{u_q}\nu_q}
 W_{d,N}(\nu_q)\otimes_{A_N}\Gamma_{d,N}(\nu_0)_r.
\end{equation}
La diferencial bar se define sin abreviación por
\begin{align}
 d_{\rm bar}(w;u_1,\ldots,u_q;c)
 :={}&(w;u_2,\ldots,u_q;\Gamma(u_1)c)
 \nonumber\\
 &+\sum_{i=1}^{q-1}(-1)^i
 (w;u_1,\ldots,u_{i+1}u_i,\ldots,u_q;c)
 \nonumber\\
 &+(-1)^q
 (W(u_q)w;u_1,\ldots,u_{q-1};c).
 \label{eq:pdfv2-coh-bar-differential}
\end{align}
Para \(q=1\), ésta es exactamente la relación
\begin{equation}\label{eq:pdfv2-coh-coend-relation}
 d_{\rm bar}(w;u;c)
 =w\otimes\Gamma(u)c-W(u)w\otimes c.
\end{equation}
Así se identifica una incidencia transportada con la misma incidencia escrita
en la carta refinada; la relación no queda oculta bajo el nombre del
totalizador.

El complejo total es
\begin{equation}\label{eq:pdfv2-coh-total-complex}
 K_{d,N}(x):=\operatorname{Tot}B_{\bullet,\bullet}^{d,N}(x),
 \qquad
 D_{\rm tot}:=d_{\rm bar}+(-1)^q\partial_\Gamma
 \quad\text{sobre }B_{q,r}^{d,N}.
\end{equation}
La functorialidad de \(W\) y \(\Gamma\), junto con
\(\partial_\Gamma^2=0\), da
\begin{equation}\label{eq:pdfv2-coh-total-square-zero}
 D_{\rm tot}^2=0.
\end{equation}
En notación compacta, pero ahora ya descargada por
\eqref{eq:pdfv2-coh-bar-bigraded}--
\eqref{eq:pdfv2-coh-bar-differential},
\begin{equation}\label{eq:pdfv2-coh-derived-coend}
 K_{d,N}(x)
 \simeq
 W_{d,N}\otimes^{\mathbf L}_{A_N[J_d(x)]}\Gamma_{d,N}.
\end{equation}

\section{Colímite de profundidad y terminal nonádico}
\label{sec:pdfv2-coh-terminal}

Los prefijos forman el colímite
\begin{equation}\label{eq:pdfv2-coh-k-infinity}
 K_{\infty,N}(x):=\varinjlim_d K_{d,N}(x).
\end{equation}
La extensión de nueve fases lleva profundidad \(d\) a profundidad \(d+9\).
Sólo después de formar \eqref{eq:pdfv2-coh-k-infinity} induce un
endomorfismo bien tipado
\begin{equation}\label{eq:pdfv2-coh-u-gamma9}
 U_{\Gamma_9,N}:K_{\infty,N}(x)
 \longrightarrow K_{\infty,N}(x).
\end{equation}
El terminal es
\begin{equation}\label{eq:pdfv2-coh-terminal-cone}
 \mathcal T_{{\rm coh},N}(x)
 :=\operatorname{Cone}\bigl(I-U_{\Gamma_9,N}\bigr).
\end{equation}
Con la convención homológica
\[
 \operatorname{Cone}(f)_q
 =K_{\infty,N,q-1}\oplus K_{\infty,N,q},
\]
su diferencial es
\begin{equation}\label{eq:pdfv2-coh-cone-differential}
 \partial_{\rm Cone}(a,b)
 =\bigl(-\partial_Ka,\partial_Kb+f(a)\bigr).
\end{equation}

Aunque la fase visible de \(h\) y \(\Gamma_9h\) coincida, los generadores
\([h]\) y \([\Gamma_9h]\) son diferentes porque sus ledgers y profundidades
son diferentes. El terminal conserva
\begin{equation}\label{eq:pdfv2-coh-terminal-difference}
 (I-U_{\Gamma_9,N})[h]=[h]-[\Gamma_9h].
\end{equation}
Las reducciones de precisión conmutan con el terminal:
\begin{equation}\label{eq:pdfv2-coh-u-reduction}
 r_{N+1,N}U_{\Gamma_9,N+1}
 =U_{\Gamma_9,N}r_{N+1,N}.
\end{equation}

\section{Lector cohomológico terminal: ciclo de incidencia y clase en \(H_0(\operatorname{Cone}(I-U_{\Gamma_9,N}))\)}
\label{sec:pdfv2-coh-reader}

Cada celda de \(x_{\leq d}\) porta un vector de incidencia
\[
 \omega_{\nu,N}(x)
 \in A_N^{I_\nu\times J_\nu}=C_{\nu,N,0}
\]
y la ruta terminal completa \(\lambda_{\nu,x}\). El ciclo bruto es
\begin{equation}\label{eq:pdfv2-coh-z-cycle}
 z_{d,N}(x):=
 \sum_{\nu\in\operatorname{Cell}_d(x)}
 \epsilon_\nu[\lambda_{\nu,x}]
 \otimes\omega_{\nu,N}(x)
 \in B_{0,0}^{d,N}(x).
\end{equation}
Como el complejo total es homológico no negativo,
\begin{equation}\label{eq:pdfv2-coh-z-closed}
 D_{\rm tot}z_{d,N}(x)=0.
\end{equation}
Se conservan tanto el ciclo como su clase
\[
 [z_{d,N}(x)]\in H_0(K_{d,N}(x)).
\]
Tras incluirlo en \(K_{\infty,N}\), su clase terminal es
\[
 [(0,z_{d,N}(x))]
 \in H_0(\mathcal T_{{\rm coh},N}(x)).
\]
El lector completo es
\begin{equation}\label{eq:pdfv2-coh-reader-full}
 \mathcal L_{{\rm coh},d,N}(x)
 :=\bigl(
 z_{d,N}(x),[z_{d,N}(x)],
 U_{\Gamma_9,N}z_{d,N}(x),
 [(0,z_{d,N}(x))]
 \bigr).
\end{equation}
No se publica únicamente una clase cociente: permanecen el representante, la
ruta que lo produjo, su continuación y el terminal.

\section{Fibras supervivientes y correcciones sin selector}
\label{sec:pdfv2-coh-surviving-fibers}

Sea \(\mathscr S_\lambda^{\rm coh}\) el \(A_N\)-módulo libre generado por las
historias finitas admisibles en el nivel \(\lambda=(d,N)\). El lector se
extiende linealmente a
\begin{equation}\label{eq:pdfv2-coh-e-lambda}
 e_\lambda:\mathscr S_\lambda^{\rm coh}
 \longrightarrow\mathscr O_\lambda^{\rm coh},
 \qquad
 \mathscr O_\lambda^{\rm coh}
 :=H_0(K_{d,N})\oplus H_0(\mathcal T_{{\rm coh},N}).
\end{equation}
El codominio superviviente es
\begin{equation}\label{eq:pdfv2-coh-surviving-output}
 \mathscr O_\infty^{\rm surv}
 :=\operatorname{Im}\!\left(
 e_\infty:\mathcal C_{\rm cont}^{\rm coh}
 \longrightarrow\varprojlim_\lambda\mathscr O_\lambda^{\rm coh}
 \right).
\end{equation}
Para \(a=(a_\lambda)_\lambda\in\mathscr O_\infty^{\rm surv}\), defínase
\begin{equation}\label{eq:pdfv2-coh-fiber}
 \mathscr F_\lambda(a):=
 \left\{s_\lambda\in\mathscr S_\lambda^{\rm coh}:
 \begin{array}{l}
 e_\lambda(s_\lambda)=a_\lambda,\\
 \exists y\in\mathcal C_{\rm cont}^{\rm coh}\text{ con }
 y|_\lambda=s_\lambda\text{ y }e_\infty(y)=a
 \end{array}
 \right\}.
\end{equation}
Estas fibras son no vacías por la definición del codominio superviviente. Si
\(\lambda'\geq\lambda\), la restricción es sobreyectiva:
\begin{equation}\label{eq:pdfv2-coh-fiber-surjection}
 \rho_{\lambda',\lambda}:\mathscr F_{\lambda'}(a)
 \twoheadrightarrow\mathscr F_\lambda(a).
\end{equation}
En efecto, una historia completa que testimonia
\(s_\lambda\in\mathscr F_\lambda(a)\) proporciona por truncación un
antecedente en cualquier nivel posterior. No se elige ese antecedente como
parte del objeto.

La multisección completa es
\begin{equation}\label{eq:pdfv2-coh-extension-multisection}
 \Sigma_{\lambda',\lambda}^{\rm coh}(s_\lambda;a)
 :=\{s_{\lambda'}\in\mathscr F_{\lambda'}(a):
 \rho_{\lambda',\lambda}s_{\lambda'}=s_\lambda\}.
\end{equation}
Para una prolongación provisional \(\widetilde s_{\lambda'}\), sea
\begin{equation}\label{eq:pdfv2-coh-discrepancy}
 \delta_{\lambda'}:=
 a_{\lambda'}-e_{\lambda'}(\widetilde s_{\lambda'}).
\end{equation}
La relación de todas las correcciones es
\begin{equation}\label{eq:pdfv2-coh-correction-relation}
 \operatorname{Corr}_{\lambda',\lambda}
 (\widetilde s_{\lambda'},\delta_{\lambda'})
 :=\left\{\eta\in\mathscr S_{\lambda'}^{\rm coh}:
 \begin{array}{l}
 \rho_{\lambda',\lambda}\eta=0,\\
 e_{\lambda'}(\eta)=\delta_{\lambda'},\\
 \eta\text{ está soportada en celdas nuevas},\\
 \eta\text{ respeta orientación, carry, memoria y terminal}
 \end{array}
 \right\}.
\end{equation}
Si dos levantamientos son comparables, su diferencia
\(\eta=s_{\lambda'}-\widetilde s_{\lambda'}\) pertenece a esta relación
exactamente cuando satisface las cuatro condiciones anteriores. No se afirma
que \eqref{eq:pdfv2-coh-correction-relation} sea no vacía para una
discrepancia arbitraria: la no vacuidad gobernante es la de
\eqref{eq:pdfv2-coh-extension-multisection} sobre el dominio superviviente.

\iffalse
\section{Las trece coordenadas de la restricción}
\label{sec:pdfv2-coh-thirteen}

Para \(x_\lambda\in\mathcal C_{{\rm cont},\lambda}^{\rm coh}\), sea
\begin{equation}\label{eq:pdfv2-coh-d-package}
 \begin{aligned}
 D_{{\rm coh},\lambda}(x_\lambda):=\bigl(&
 \mathcal F_{{\rm coh},\lambda},
 \operatorname{Ext}_{\Gamma,{\rm coh},\lambda},
 \operatorname{Rel}_{{\rm coh},\lambda},
 \mathcal N_{{\rm coh},\lambda},
 \operatorname{or}_{{\rm coh},\lambda},
 \operatorname{Carr}_{{\rm coh},\lambda},
 \operatorname{Mem}_{{\rm coh},\lambda},\\
 &\partial_{{\rm coh},\lambda},
 \gamma_{{\rm coh},\lambda},
 \operatorname{Msect}_{{\rm coh},\lambda},
 \operatorname{Surv}_{{\rm coh},\lambda},
 \operatorname{Term}_{{\rm coh},\lambda},
 \mathcal L_{{\rm coh},\lambda}\bigr).
 \end{aligned}
\end{equation}
Sus coordenadas son:
\begin{enumerate}
 \item \(\mathcal F_{{\rm coh},\lambda}\): todos los módulos
       \(A_N^{I_\nu}\), \(A_N^{J_\nu}\),
       \(A_N^{I_\nu\times J_\nu}\), los complejos \(C_{\nu,N}\), los
       módulos \(W_{d,N}(\nu)\) y \(\mathscr S_\lambda^{\rm coh}\).
 \item \(\operatorname{Ext}_{\Gamma,{\rm coh},\lambda}\): todos los mapas
       \(\Gamma(u)\), reducciones de coeficientes, truncaciones, extensiones
       \(\Gamma_9\) y las relaciones \(\Sigma_{\lambda',\lambda}\).
 \item \(\operatorname{Rel}_{{\rm coh},\lambda}\):
       \(\kappa\), \(\delta\), la sucesión exacta, las caras bar, la
       relación de coend, la orientación y las leyes de composición.
 \item \(\mathcal N_{{\rm coh},\lambda}\):
       \[
       (d,N,3^N,a_\nu,b_\nu,(a_\nu-1)(b_\nu-1),
       \rho_{\Sigma,\nu},\rho_{\Pi,\nu},c_{\Sigma,\nu},c_{\Pi,\nu}).
       \]
 \item \(\operatorname{or}_{{\rm coh},\lambda}\): los signos
       \(\epsilon_\nu\), los órdenes de puertos y los mapas
       \(\Theta_{\nu,N}\).
 \item \(\operatorname{Carr}_{{\rm coh},\lambda}\): las matrices
       \(\beta_N\), los cociclos \(c_N(v,u)\), residuos y cocientes,
       conservados por separado.
 \item \(\operatorname{Mem}_{{\rm coh},\lambda}\): el ledger heredado y la
       palabra ordenada de refinamientos, carries, correcciones y avances
       \(\Gamma_9\), no sólo su valor agregado.
 \item \(\partial_{{\rm coh},\lambda}\):
       \(\kappa_{\nu,N}\), \(D_{\rm tot}\),
       \(\partial_{\rm Cone}\) y las fronteras de las celdas.
 \item \(\gamma_{{\rm coh},\lambda}\): todas las cadenas componibles
       \(\nu_0\to\cdots\to\nu_q\) y sus continuaciones terminales.
 \item \(\operatorname{Msect}_{{\rm coh},\lambda}\): todas las celdas,
       rutas terminales, extensiones \(\Sigma\) y correcciones
       \(\operatorname{Corr}\), sin selector.
 \item \(\operatorname{Surv}_{{\rm coh},\lambda}\): el sistema de fibras
       \(\mathscr F_\lambda(a)\) y restricciones sobreyectivas.
 \item \(\operatorname{Term}_{{\rm coh},\lambda}\):
       \(\operatorname{Cone}(I-U_{\Gamma_9,N})\), junto con \(z\), \(Uz\)
       y su diferencia.
 \item \(\mathcal L_{{\rm coh},\lambda}\):
       \((z,[z],Uz,[(0,z)])\) y todas sus compatibilidades de reducción.
\end{enumerate}

Para \(\lambda'\geq\lambda\), el transporte coordinado satisface
\begin{equation}\label{eq:pdfv2-coh-d-naturality}
 u_{\lambda',\lambda}^{\rm coh}
 D_{{\rm coh},\lambda'}(x_{\lambda'})
 =D_{{\rm coh},\lambda}
  (\rho_{\lambda',\lambda}x_{\lambda'}).
\end{equation}

\section{Grafo de restricción y objeto dependiente}
\label{sec:pdfv2-coh-graphs}

La restricción total es el grafo
\begin{equation}\label{eq:pdfv2-coh-g-graph}
 \mathcal G_{\rm coh}:=\operatorname{Graph}(D_{\rm coh}),
 \qquad
 \operatorname{Res}_{\rm coh}(x):=(x,D_{\rm coh}x).
\end{equation}
Así,
\begin{equation}\label{eq:pdfv2-coh-res-inverse}
 \pi_1\operatorname{Res}_{\rm coh}=I,
 \qquad
 \operatorname{Res}_{\rm coh}\pi_1=I_{\mathcal G_{\rm coh}}.
\end{equation}

El extractor completo es
\begin{equation}\label{eq:pdfv2-coh-chi}
 \chi_{\rm coh}(x):=
 \bigl(
 \{C_{\nu,N}\},
 \{z_\lambda,[z_\lambda]\},
 \{\mathscr F_\lambda(e_\infty x)\},
 \{\Sigma_{\lambda',\lambda}\},
 \{\operatorname{Corr}_{\lambda',\lambda}\}
 \bigr).
\end{equation}
Se define
\begin{equation}\label{eq:pdfv2-coh-xchi}
 X_{\chi,{\rm coh}}
 :=\{(\chi_{\rm coh}(x),x,D_{\rm coh}x):
 x\in\mathcal C_{\rm cont}^{\rm coh}\}
\end{equation}
y
\begin{equation}\label{eq:pdfv2-coh-prx}
 \operatorname{pr}_{X,{\rm coh}}(x,D_{\rm coh}x)
 :=(\chi_{\rm coh}(x),x,D_{\rm coh}x).
\end{equation}
La proyección sobre \((x,D_{\rm coh}x)\) es inversa sobre este objeto
dependiente.

\section{Ensamblador y publicador internos}
\label{sec:pdfv2-coh-assembler-publisher}

El ambiente \(\mathscr M_{\rm coh}^{\rm amb}\) consta de sistemas
compatibles de sucesiones exactas de celdas, categorías, funtores, módulos de
ruta, totalizadores, terminales, lectores, fibras y multisecciones. El
ensamblador bruto es
\begin{equation}\label{eq:pdfv2-coh-assembler}
 \begin{aligned}
 a_{\rm coh}:X_{\chi,{\rm coh}}&\longrightarrow
 \mathscr M_{\rm coh}^{\rm amb},\\
 a_{\rm coh}(\chi(x),x,Dx)
 &:={}
 \bigl(
 \{C_{\nu,N},\kappa_{\nu,N},\delta_{\nu,N}\},
 \{J_d,\Gamma,W,K\},\\
 &\hspace{31mm}
 \{U_{\Gamma_9,N},\mathcal T_N\},
 \{e_\lambda,\mathscr F_\lambda,\Sigma,\operatorname{Corr}\}
 \bigr)_x.
 \end{aligned}
\end{equation}
La macroestructura es otro grafo:
\begin{equation}\label{eq:pdfv2-coh-macro-graph}
 \mathfrak M_{\rm coh}:=\operatorname{Graph}(a_{\rm coh}),
 \qquad
 \mathcal A_{\rm coh}(z):=(z,a_{\rm coh}z).
\end{equation}

Sea \(\mathscr Y_{\rm coh}^{\rm int}\) el tipo de certificados que contiene
\begin{equation}\label{eq:pdfv2-coh-certificate-list}
 \begin{gathered}
 \text{la sucesión exacta \eqref{eq:pdfv2-coh-exact-sequence}},\\
 \kappa_{b,a}\tau=-P\kappa_{a,b},
 \qquad D_{\rm tot}^2=0,
 \qquad\partial_{\rm Cone}^2=0,\\
 rU_{\Gamma_9}=U_{\Gamma_9}r,
 \qquad
 e_\lambda\rho_{\lambda',\lambda}
 =r_{\lambda',\lambda}e_{\lambda'},\\
 \rho_{\lambda',\lambda}:\mathscr F_{\lambda'}(a)
 \twoheadrightarrow\mathscr F_\lambda(a),
 \qquad
 \Sigma_{\lambda',\lambda}(s_\lambda;a)\neq\varnothing,
 \end{gathered}
\end{equation}
junto con los ciclos brutos, sus clases y sus terminales. El publicador
\begin{equation}\label{eq:pdfv2-coh-publisher}
 p_{\rm coh}^{\rm raw}:\mathfrak M_{\rm coh}
 \longrightarrow\mathscr Y_{\rm coh}^{\rm int}
\end{equation}
produce ese certificado sin eliminar la macroestructura. Finalmente,
\begin{equation}\label{eq:pdfv2-coh-output-graph}
 \mathcal C_{\rm coh}^{\rm HMT}
 :=\operatorname{Graph}(p_{\rm coh}^{\rm raw}),
 \qquad
 \mathcal P_{\rm coh}(m):=(m,p_{\rm coh}^{\rm raw}m).
\end{equation}

\section{Teorema interno de generación coherente}
\label{sec:pdfv2-coh-theorem}

\begin{theorem}[Generación interna de la macroestructura de coherencia]
\label{thm:pdfv2-coh-internal-generation}
Sobre el dominio explícito \(\mathcal C_{\rm cont}^{\rm coh}\), la cadena
\[
 \mathcal C_{\rm cont}^{\rm coh}
 \xrightarrow{\operatorname{Res}_{\rm coh}}
 \mathcal G_{\rm coh}
 \xrightarrow{\operatorname{pr}_{X,{\rm coh}}}
 X_{\chi,{\rm coh}}
 \xrightarrow{\mathcal A_{\rm coh}}
 \mathfrak M_{\rm coh}
 \xrightarrow{\mathcal P_{\rm coh}}
 \mathcal C_{\rm coh}^{\rm HMT}
\]
construye de forma natural las celdas, sus conos, el diagrama de historias, el
bar de dos lados, la totalización, el terminal nonádico y las fibras
supervivientes. Conserva las trece coordenadas y todas las correcciones como
multisecciones. Cada etapa admite, sobre su imagen, una proyección que recupera
la etapa anterior.
\end{theorem}

\begin{proof}
La \cref{prop:pdfv2-coh-cell-exact} construye cada celda sobre \(A_N\), con
núcleo diagonal y defecto \((a_\nu-1)(b_\nu-1)\). La identidad
\eqref{eq:pdfv2-coh-orientation} hace del intercambio orientado un mapa de
complejos involutivo. Las fórmulas
\eqref{eq:pdfv2-coh-gamma-one}--\eqref{eq:pdfv2-coh-gamma-zero} prueban que
cada refinamiento es un mapa de complejos, mientras que
\eqref{eq:pdfv2-coh-kappa-reduction} prueba la compatibilidad en precisión.

La precomposición de rutas hace contravariante a \(W\). Las identidades
functoriales de \(W\) y \(\Gamma\) implican
\(d_{\rm bar}^2=0\); su conmutación con la diferencial celular da
\(D_{\rm tot}^2=0\). Así, \(K_{d,N}\) queda construido por las fórmulas bar,
no por una invocación nominal del coend.

La extensión \(\Gamma_9\) aumenta la profundidad. Por ello sólo induce el
endomorfismo \(U_{\Gamma_9,N}\) después de pasar al colímite
\(K_{\infty,N}\). El cono
\(\operatorname{Cone}(I-U_{\Gamma_9,N})\) conserva entonces la diferencia
\([h]-[\Gamma_9h]\) y no identifica la repetición de fase con la repetición
de estado.

El ciclo \(z_{d,N}(x)\) se forma con todas las celdas, signos, rutas e
incidencias de la historia. Truncación y reducción conmutan con el lector. Si
\(s_\lambda\in\mathscr F_\lambda(a)\), la historia completa que aparece en
la propia definición de la fibra proporciona un antecedente en todo nivel
\(\lambda'\geq\lambda\); esto prueba la sobreyectividad
\eqref{eq:pdfv2-coh-fiber-surjection}. El objeto conserva el conjunto entero
\(\Sigma_{\lambda',\lambda}\), no una elección. Las diferencias entre
levantamientos se registran mediante
\eqref{eq:pdfv2-coh-correction-relation}, con soporte, orientación, carry,
memoria y terminal explícitos.

Finalmente, \(\mathcal G_{\rm coh}\), \(\mathfrak M_{\rm coh}\) y
\(\mathcal C_{\rm coh}^{\rm HMT}\) son grafos, mientras que
\(X_{\chi,{\rm coh}}\) conserva como coordenadas \(x\) y \(D_{\rm coh}x\).
La primera proyección de cada objeto recupera el anterior. Por tanto, la
composición completa no sustituye historia por clase, multisección por
selector ni terminal por fase visible.
\end{proof}

\section{Procedencia probatoria y falsadores}
\label{sec:pdfv2-coh-provenance-falsifiers}

\paragraph{Resultado recuperado.}
Pertenecen a la arquitectura ya construida la celda \(\kappa_{a,b}\), su
núcleo diagonal y defecto, la orientación
\(\kappa_{b,a}\tau=-P\kappa_{a,b}\), el carry APP, el cono de la celda, la
categoría de historias, la totalización por coend, la acción \(\Gamma_9\), el
terminal \(\operatorname{Cone}(I-U_{\Gamma_9})\) y la arquitectura de
extensión superviviente.

\paragraph{Formalización nueva.}
Son formalizaciones explícitas la sucesión exacta
\eqref{eq:pdfv2-coh-exact-sequence} sobre \(A_N\), la separación entre
profundidad y precisión, las caras completas del bar, la formación del
terminal sólo después del colímite, el empaquetado de las trece coordenadas,
las correcciones como relación sin selector y los cuatro objetos dependientes
o grafos que impiden pérdida de genealogía.

\begin{criterion}[Falsadores internos de la rama \(\mathrm{coh}\)]
\label{crit:pdfv2-coh-falsifiers}
La construcción falla si ocurre cualquiera de los hechos siguientes:
\begin{enumerate}
 \item no es exacta la sucesión
       \eqref{eq:pdfv2-coh-exact-sequence};
 \item falla \(\kappa_{b,a}\tau=-P\kappa_{a,b}\);
 \item algún refinamiento no satisface
       \(\kappa_{\nu'}\Gamma(u)_1=\Gamma(u)_0\kappa_\nu\);
 \item el carry de composición no cumple
       \eqref{eq:pdfv2-coh-carry-cocycle};
 \item \(D_{\rm tot}^2\neq0\);
 \item reducción, lector o terminal no conmutan con \(U_{\Gamma_9}\);
 \item se forma \(\operatorname{Cone}(I-U_{\Gamma_9})\) identificando antes
       las profundidades \(d\) y \(d+9\);
 \item una fibra declarada superviviente es vacía o una restricción no es
       sobreyectiva;
 \item se sustituye \(\Sigma_{\lambda',\lambda}\) por un único
       levantamiento;
 \item se publica sólo \([z]\) y se eliminan \(z\), ruta, carry, memoria o
       terminal;
 \item se omite cualquiera de las trece coordenadas.
\end{enumerate}
En el último caso el diagnóstico vinculante es: \emph{objeto sustituido;
conclusión no transferible}.
\end{criterion}

\begin{warningbox}
Este capítulo concluye únicamente la construcción interna de
\(\mathcal C_{\rm coh}^{\rm HMT}\). No introduce un espacio clásico, una
clase clásica ni una publicación clásica. Cualquier entrelazador posterior
pertenece a otro documento y no gobierna la existencia de esta
macroestructura.
\end{warningbox}
\fi
